% =====================================================================
% From the DSGE state-space solution to a VARMA and to a VAR(infinity)
% Derivation in the notation of Ravenna (2007, JME 54, 2048--2064)
% Cross-referenced to Fernandez-Villaverde, Rubio-Ramirez & Sargent
%   ("A, B, C's (and D)'s for Understanding VARs", NBER TWP 308, 2005)
% =====================================================================

\section{From the DSGE solution to a VAR($\infty$): the truncation problem}
\label{sec:varma}

This section establishes the theoretical mapping between the state-space
solution of a linearised DSGE model, its VARMA reduced form, and the
implied infinite-order VAR. The purpose is to show that a finite-order
VAR($p$) is, generically, a \emph{misspecified} representation of the
model's data-generating process, and that the severity of the resulting
truncation bias is governed by an object that can be computed directly
from the model solution.

\subsection{The state-space solution}

Following \citet{Ravenna2007}, the recursive equilibrium law of motion of
a linearised DSGE model can be written as
\begin{align}
    y_t &= P x_{t-1} + Q z_t, \label{eq:rav1}\\
    x_t &= R x_{t-1} + S z_t, \label{eq:rav2}\\
    Z(L) z_t &= \varepsilon_t, \label{eq:rav3}
\end{align}
where $x_t$ is the $n \times 1$ vector of \emph{endogenous} state
variables, $z_t$ is the $m \times 1$ vector of \emph{exogenous} state
variables (the structural shocks), $y_t$ is the $r \times 1$ vector of
endogenous control variables, and $\varepsilon_t$ is an $m \times 1$
white-noise process with $\mathbb{E}(\varepsilon_t) = 0$,
$\mathbb{E}(\varepsilon_t \varepsilon_t') = \Sigma$ diagonal. The matrix
polynomial $Z(L) = [I - Z_1 L - \dots - Z_q L^q]$ defines a stationary
$\mathrm{VAR}(q)$ process for the exogenous states.

The system \eqref{eq:rav1}--\eqref{eq:rav3} is the exact analogue of the
$(A,B,C,D)$ representation of \citet{FVRRS2005}: their state equation
$X_{t+1} = A X_t + B W_{t+1}$ corresponds to \eqref{eq:rav2} and their
observation equation $Y_t = C X_t + D W_t$ to \eqref{eq:rav1}. Ravenna's
partition into endogenous states $x_t$ and controls $y_t$ is more
convenient here, because the two objects play different roles in what
follows.

\subsection{Step 1: eliminating the endogenous states}

Write \eqref{eq:rav2} in lag-operator form as $G(L)\, x_t = S z_t$, with
\begin{equation}
    G(L) \equiv [I - R L],
\end{equation}
and assume $G(L)$ is invertible, so that
$x_{t-1} = G(L)^{-1} S L\, z_t$. Substituting into \eqref{eq:rav1} and
taking as the observable vector $\widehat{Y}_t = y_t$ --- that is,
assuming the econometrician observes the controls but \emph{not} the
endogenous states --- yields
\begin{equation}
    y_t = Q z_t + P\, G(L)^{-1} S L\, z_t.
    \label{eq:rav21}
\end{equation}
This is Ravenna's equation (21): an infinite-order vector moving average
in the structural shocks. \textbf{The moving-average term arises precisely
because $x_t$ is not observed.} If all states were included in the
observable vector, the term would collapse and a finite-order VAR would
obtain directly.

\subsection{Step 2: the VARMA representation (i.i.d. shocks)}

Consider first the simplest case $Z(L) = I$, so that $z_t =
\varepsilon_t$. Assume $Q$ is invertible --- which requires $r = m$, i.e.
\emph{as many observables as structural shocks} --- and define the
reduced-form innovation vector
\begin{equation}
    \eta_t \equiv Q \varepsilon_t = Q z_t .
\end{equation}
The matrix $Q$ is thus the theoretical identifying matrix mapping
structural shocks into reduced-form innovations. Equation
\eqref{eq:rav21} becomes
$y_t = \eta_t + P G(L)^{-1} S Q^{-1} L\, \eta_t$.

To clear the inverse, use the determinant--adjoint decomposition
$G(L)^{-1} = D_G(L)\,|G(L)|^{-1}$, where the scalar polynomial $|G(L)|$
is of order $n$ in $L$ and the adjoint matrix $D_G(L)$ is of order $n-1$.
Multiplying through by $|G(L)|$:
\begin{equation}
    |G(L)|\, y_t = |G(L)|\, \eta_t + P D_G(L) S Q^{-1} L\, \eta_t
    \;\equiv\; G^{*}(L)\, \eta_t .
    \label{eq:rav22}
\end{equation}
Equation \eqref{eq:rav22} --- Ravenna's equation (22) --- is a
$\mathrm{VARMA}(n,n)$ in \emph{final equations form}: the autoregressive
part is the scalar polynomial $|G(L)|$ times the identity, so each
component of $y_t$ depends only on its own lags. Since the zero-order
coefficient of both $|G(L)|$ and $G^{*}(L)$ is the identity, this
representation is unique.

\subsection{Step 3: from VARMA to VAR($\infty$)}

If $G^{*}(L)$ is invertible, inverting the moving-average part of
\eqref{eq:rav22} delivers the autoregressive representation
\begin{equation}
    |G(L)|\, G^{*}(L)^{-1} y_t = \eta_t .
    \label{eq:rav23}
\end{equation}
Writing $G^{*}(L)^{-1} = D_{G^{*}}(L)\,|G^{*}(L)|^{-1}$, the adjoint
$D_{G^{*}}(L)$ is a finite-order lag polynomial, but the inverse of the
scalar operator $|G^{*}(L)|$ is of \emph{infinite} order unless
$|G^{*}(L)|$ is constant. Hence:

\begin{quote}
\textbf{Generically, the DSGE solution implies a VAR of infinite order.}
A finite-order VAR($p$) representation exists if and only if
$|G^{*}(L)|$ is of degree zero in $L$, i.e. $G^{*}(L)$ is unimodular
\citep[p.~245]{Lutkepohl1993}.
\end{quote}

\subsection{Step 4: the persistent-shock case}

The empirically relevant case has $Z(L) \neq I$ (autocorrelated
structural shocks, as in ACEL). Define
$\widetilde{Z}(L) \equiv Q Z(L) Q^{-1} =
[I - Q Z_1 Q^{-1} L - \dots - Q Z_q Q^{-1} L^q]$, so that
$Q Z(L)^{-1}\varepsilon_t = \widetilde{Z}(L)^{-1}\eta_t$. Then
\begin{equation}
    |\widetilde{Z}(L)|\,|G(L)|\, y_t
    = G^{*}(L)\, D_{\widetilde{Z}}(L)\, \eta_t ,
    \label{eq:rav24}
\end{equation}
which, since $|\widetilde{Z}(L)|$ is of order $qm$ in $L$, is a
\begin{equation}
    \mathrm{VARMA}\big(n + qm,\; n + q(m-1)\big)
\end{equation}
process --- Proposition 2.1 of \citet{Ravenna2007}. The corresponding
autoregressive representation,
$\widetilde{Z}(L)\,G^{*}(L)^{-1}|G(L)|\,y_t = \eta_t$, is again of
infinite order unless the unimodularity condition above is satisfied.

\subsection{Step 5: the decay rate of the VAR($\infty$) coefficients}

Corollary 2.2 of \citet{Ravenna2007} restates the finite-order condition
in a form that is directly computable and that links back to
\citet{FVRRS2005}. From \eqref{eq:rav1}, $z_{t-1} = Q^{-1}(y_{t-1} - P
x_{t-1})$; substituting into \eqref{eq:rav2} gives $H(L)\,x_t = S Q^{-1}
y_{t-1}$ with
\begin{equation}
    H(L) \equiv \big[\,I - (R - S Q^{-1} P)\,L\,\big].
    \label{eq:H}
\end{equation}
If $H(L)$ is invertible, the observables admit the autoregressive
representation
\begin{equation}
    y_t
    = P \sum_{j=0}^{\infty} \big[\,R - S Q^{-1} P\,\big]^{\,j}
      S Q^{-1} L^{\,j+1} y_t \;+\; \eta_t .
    \label{eq:rav26}
\end{equation}

Three observations follow, and they are central to this thesis.

\paragraph{(i) The bridge to the ABCD condition.} The matrix
$R - S Q^{-1} P$ in \eqref{eq:H} is the exact counterpart of
$A - B D^{-1} C$ in \citet{FVRRS2005}. Their ``poor man's invertibility
condition'' --- all eigenvalues of $A - BD^{-1}C$ strictly inside the
unit circle --- is therefore the same requirement that makes the sum in
\eqref{eq:rav26} converge. Ravenna notes explicitly that
\eqref{eq:rav26} ``is the derivation obtained in Fernandez-Villaverde et
al.\ (2005)'' \citep{FVRRS2005}.

\paragraph{(ii) Existence versus accuracy.} The eigenvalue condition
guarantees that the VAR($\infty$) \emph{exists}. It says nothing about
how fast the coefficient matrices
$P\,[R - SQ^{-1}P]^{\,j}\,SQ^{-1}$ decay in $j$. The decay is governed by
the \emph{modulus of the largest eigenvalue} of $R - SQ^{-1}P$: the
closer it lies to unity, the more slowly the coefficients vanish, and the
more severe the bias induced by truncating the lag structure at some
finite $p$. This is the mechanism behind the finding of
\citet{ChariKehoeMcGrattan2005} that the exact VAR($\infty$) coefficients
of a standard business-cycle model ``converge to zero extremely slowly'',
making short-lag VAR approximations unsuitable.

\paragraph{(iii) Implication for the exercise in this thesis.} Truncation
bias is therefore not a small-sample phenomenon and does not depend on a
mistaken identification scheme: it is a population property of the mapping
between the DSGE and the finite-order VAR. \citet{Ravenna2007} shows that
identification strategies which are correct in the exact VAR($\infty$)
can perform poorly once the lag structure is truncated. This provides the
theoretical justification for enriching the lag structure of the
empirical VAR, which is the exercise conducted in the remainder of this
thesis.

\subsection{Stochastic singularity and the status of the ACEL VAR}
\label{sec:singularity}

The derivation above requires $Q$ to be invertible, that is, the number
of observables entering the system must equal the number of structural
shocks, $r = m$. The ACEL framework does not satisfy this condition: the
model is driven by $m = 3$ structural shocks (monetary policy, neutral
technology, embodied technology) while the empirical VAR contains
$10$ observable series. In Ravenna's terminology this is the case
$m < n + r$, for which he notes that ``the system is singular, preventing
likelihood estimation of the VAR'' \citep[p.~2050]{Ravenna2007}: a vector
of ten series driven by three fundamental shocks has a reduced-form
innovation covariance matrix of rank at most three, which cannot be
inverted.

Two consequences follow, and it is worth stating them explicitly because
they delimit what the theory above does and does not deliver.

First, \emph{the singularity is not resolved by moving to a VARMA
representation}. The moving-average terms in \eqref{eq:rav22} and
\eqref{eq:rav24} arise from the omission of the endogenous \emph{states}
$x_t$ from the observable vector, not from an excess of observables over
shocks. A VARMA for the full ten-variable vector would inherit the same
rank-deficient innovation covariance. Recovering an exact VARMA, and
hence an exact VAR($\infty$), requires selecting a sub-vector of
observables of dimension equal to $m$, as in the applied strategy of
\citet{CEV2006}.

Second, this is precisely why the empirical strategy of
\citet{Altigetal2011} --- and hence of this thesis --- proceeds by
\emph{minimum-distance impulse-response matching} rather than by
full-information likelihood estimation of a joint DSGE--VAR system. The
ten-variable VAR is best understood as a flexible reduced-form
approximation to the dynamics of the data, against which the model's
impulse responses are evaluated, rather than as the exact
data-generating process implied by the model. The truncation-bias results
derived above then apply to this VAR as an approximation property: even
abstracting from singularity, the model-consistent representation of any
subset of observables is generically of infinite order, so the choice of
$p$ is a substantive specification decision rather than an innocuous
one.

% ---------------------------------------------------------------------
% Suggested bibliography keys (add to biblio.bib if not present):
%   Ravenna2007, FVRRS2005, Lutkepohl1993,
%   ChariKehoeMcGrattan2005, CEV2006, Altigetal2011
% ---------------------------------------------------------------------
